Terms, acronyms and platform-specific identifiers used in the paper and
its appendices. Grouped by domain.

\begin{center}\rule{0.5\linewidth}{0.5pt}\end{center}

\subsection{Embedded-systems / RTOS}\label{embedded-systems-rtos}

{\def\LTcaptype{none} % do not increment counter
\begin{longtable}[]{@{}
  >{\raggedright\arraybackslash}p{(\linewidth - 2\tabcolsep) * \real{0.5000}}
  >{\raggedright\arraybackslash}p{(\linewidth - 2\tabcolsep) * \real{0.5000}}@{}}
\toprule\noalign{}
\begin{minipage}[b]{\linewidth}\raggedright
Term
\end{minipage} & \begin{minipage}[b]{\linewidth}\raggedright
Expansion / meaning
\end{minipage} \\
\midrule\noalign{}
\endhead
\bottomrule\noalign{}
\endlastfoot
\textbf{DMA} & Direct Memory Access. Hardware engine that moves bytes
without CPU involvement. \\
\textbf{eDMA} & ``enhanced'' DMA --- the NXP i.MX RT1176 DMA0 peripheral
used for the SDRAM→tensor memcpy. \\
\textbf{ISR} & Interrupt Service Routine. Runs in interrupt context with
bounded, minimal work per NASA/JPL discipline
(\href{../agent/embeded.md}{embeded.md} §C). \\
\textbf{MCU} & Microcontroller Unit. Here: NXP i.MX RT1176. \\
\textbf{RTOS} & Real-Time Operating System. Here: FreeRTOS 10.x (CMSIS
M7 build). \\
\textbf{SDRAM} & Synchronous DRAM. On-module, 16 MB on the SEMC bus at
166 MHz. \\
\textbf{SEMC} & Smart External Memory Controller. The RT1176 peripheral
that drives the external SDRAM. \\
\textbf{SDP} & Serial Download Protocol. NXP's ROM bootloader protocol;
the board falls back to SDP if firmware fails to boot. \\
\textbf{WDOG} & Hardware watchdog (WDOG1). 30 s timeout on the 32 kHz
clock, independent of CPU. \\
\textbf{VBLANK} & Vertical blanking interval. Period between the last
line of one video frame and the first line of the next, during which no
pixel data flows on the MIPI-CSI lane. \\
\textbf{VSYNC / HSYNC} & Vertical / horizontal sync pulses marking frame
/ line boundaries. In MIPI-CSI2, replaced by FS / LS short packets. \\
\end{longtable}
}

\begin{center}\rule{0.5\linewidth}{0.5pt}\end{center}

\subsection{Camera / imaging}\label{camera-imaging}

{\def\LTcaptype{none} % do not increment counter
\begin{longtable}[]{@{}
  >{\raggedright\arraybackslash}p{(\linewidth - 2\tabcolsep) * \real{0.5000}}
  >{\raggedright\arraybackslash}p{(\linewidth - 2\tabcolsep) * \real{0.5000}}@{}}
\toprule\noalign{}
\begin{minipage}[b]{\linewidth}\raggedright
Term
\end{minipage} & \begin{minipage}[b]{\linewidth}\raggedright
Meaning
\end{minipage} \\
\midrule\noalign{}
\endhead
\bottomrule\noalign{}
\endlastfoot
\textbf{CSI} & Camera Serial Interface. The RT1176 on-chip peripheral
that reads pixel data. \\
\textbf{CSI-2} & MIPI Camera Serial Interface, version 2. The protocol
spoken between the OV5640 and the RT1176 on our platform. \\
\textbf{D-PHY} & MIPI physical-layer specification that CSI-2 runs
over. \\
\textbf{FS / LS} & Frame Start / Line Start --- short packets in
MIPI-CSI2 marking frame and line boundaries. \\
\textbf{EOF} & End-of-Frame --- the interrupt fired when a DMA buffer
has finished filling. Our flip-on-EOF fix runs inside this ISR. \\
\textbf{MUX} & Multiplexer. On this board, an analogue switch on the
shared MIPI-CSI2 lane, GPIO-controlled; selects which of the two OV5640
sensors the CSI receiver is talking to. \\
\textbf{PXP} & Pixel Pipeline. The RT1176 2-D graphics accelerator used
for colour conversion and resizing between the raw DMA buffer and the
TPU input tensor. \\
\textbf{OV5640} & OmniVision 5-megapixel image sensor used as both cam0
and cam1 on this board. Dual instance, shared MIPI lane via MUX. \\
\textbf{FSIN} & Frame Sync Input pin on the OV5640, used for
master/slave synchronisation between two sensors. Not wired on the
SentAI board. \\
\textbf{AEC / AGC} & Auto Exposure Control / Auto Gain Control --- the
sensor's internal adaptive exposure loops. Take several frames to
converge after a stream resume. \\
\textbf{QSXGA / 1080P / 720P / VGA / QVGA} & Standard resolutions:
2592×1944, 1920×1080, 1280×720, 640×480, 320×240. \\
\end{longtable}
}

\begin{center}\rule{0.5\linewidth}{0.5pt}\end{center}

\subsection{AI inference}\label{ai-inference}

{\def\LTcaptype{none} % do not increment counter
\begin{longtable}[]{@{}
  >{\raggedright\arraybackslash}p{(\linewidth - 2\tabcolsep) * \real{0.5000}}
  >{\raggedright\arraybackslash}p{(\linewidth - 2\tabcolsep) * \real{0.5000}}@{}}
\toprule\noalign{}
\begin{minipage}[b]{\linewidth}\raggedright
Term
\end{minipage} & \begin{minipage}[b]{\linewidth}\raggedright
Meaning
\end{minipage} \\
\midrule\noalign{}
\endhead
\bottomrule\noalign{}
\endlastfoot
\textbf{EdgeTPU} & Google-designed systolic-array AI accelerator, on-die
on the Coral Dev Board Micro. Runs 8-bit quantised TFLite models at
fixed clock; execution time scales with input pixel count and FLOPs. \\
\textbf{TFLite} & TensorFlow Lite. The on-device inference framework.
``TFLite Micro'' is the embedded variant we use via
\texttt{sentai.tpu.*}. \\
\textbf{YOLOv5 / YOLOv5-enhanced} & Single-stage object-detection
architecture family. The ``enhanced'' 1-class 512×512 model used
throughout E15-E18 has a single upsample at P5/32 and a 1-class head
producing \texttt{{[}1,\ 1344,\ 6{]}}. \\
\textbf{NMS} & Non-Maximum Suppression. Post-processing step that
deduplicates overlapping detections. Runs on the CPU
(MicroPython-callable) after Invoke. \\
\textbf{Anchor} & A pre-defined bounding-box prior used by YOLO-style
detectors. 1344 anchors in the 512×512 single-upsample model. \\
\textbf{Mode 3 (\texttt{kMax})} & EdgeTPU performance mode requested by
\texttt{sentai.tpu.load()}. \\
\textbf{Invoke} & Single forward pass of the TFLite model. Reported as
\texttt{invoke\_ms} in every CSV. \\
\textbf{Quantisation (scale, zero\_point)} & 8-bit integer
representation of activations;
\texttt{real\ =\ (int\ −\ zero\_point)\ ×\ scale}. \\
\end{longtable}
}

\begin{center}\rule{0.5\linewidth}{0.5pt}\end{center}

\subsection{Firmware platform
(SentAI-specific)}\label{firmware-platform-sentai-specific}

{\def\LTcaptype{none} % do not increment counter
\begin{longtable}[]{@{}
  >{\raggedright\arraybackslash}p{(\linewidth - 2\tabcolsep) * \real{0.5000}}
  >{\raggedright\arraybackslash}p{(\linewidth - 2\tabcolsep) * \real{0.5000}}@{}}
\toprule\noalign{}
\begin{minipage}[b]{\linewidth}\raggedright
Term
\end{minipage} & \begin{minipage}[b]{\linewidth}\raggedright
Meaning / where defined
\end{minipage} \\
\midrule\noalign{}
\endhead
\bottomrule\noalign{}
\endlastfoot
\textbf{\texttt{sentai} module} & Root MicroPython namespace that
exposes the hardware surface. Submodules: \texttt{camera}, \texttt{tpu},
\texttt{pipeline}, \texttt{fs}, \texttt{rtos}, \texttt{diag},
\texttt{io}, \texttt{imu}, \texttt{mic}, \texttt{usb}, \texttt{link},
\texttt{mesh}, \texttt{crazy}, \texttt{tfl}, \texttt{aifes},
\texttt{kmeans}, \texttt{pca}, \texttt{anomaly}, \texttt{dtw},
\texttt{hmm}, \texttt{rl}, \texttt{slam}. \\
\textbf{PrepTask / InferTask} & The two tasks inside the firmware
parallel pipeline (see \url{camera.md} §PXP). PrepTask does grab + PXP +
quant into a staging buffer; InferTask does memcpy-to-tensor + Invoke +
NMS. \\
\textbf{Staging buffer / tensor buffer} & The two 512×512×3 buffers
between the two pipeline tasks. The staging buffer is written by
PrepTask; the tensor buffer is written by InferTask's memcpy (eDMA since
the fix) and read by the TPU. \\
\textbf{\texttt{sentai.camera.select(id)}} & Arm a glitch-free flip to
camera \texttt{id}; the CSI EOF ISR consumes the arm in VBLANK.
\url{cam_switch.md} §``Fix B''. \\
\textbf{\texttt{sentai.camera.switch\_drain(n)}} & Number of fresh ISR
frames required after a MUX flip before the next captured frame is
returned. Default 2; 1 is an experimentation hook. \\
\textbf{\texttt{sentai.camera.ratio(a,\ b)}} & Stateless auto-alternate
scheduler: over any \texttt{(a+b)}-frame cycle, cam0 gets \texttt{a}
frames and cam1 gets \texttt{b}. \texttt{(0,\ 0)} disables. \\
\textbf{\texttt{sentai.diag.cam\_stats()}} & Returns a dict of
persistent fault counters: \texttt{switch\_ok\_eof},
\texttt{switch\_fallback}, \texttt{drain\_timeout},
\texttt{grab\_retry}, \texttt{grab\_fatal}. \\
\textbf{Session} & A
\texttt{/diags/sNNN\_\textless{}name\textgreater{}/} folder on the
LittleFS user partition, opened by \texttt{diag.begin(...)} and closed
by \texttt{diag.end()}. Self-contained: CSV + descriptor text + manifest
+ summary + scene snapshots. \\
\textbf{Manifest} & \texttt{manifest.csv} inside a session: one row per
experiment completed in that session. \\
\textbf{Diag experiment} & A function
\texttt{e\textless{}N\textgreater{}\_xxx(...)} inside the \texttt{diag}
package; the numbered experiment classes referenced throughout this
paper (E13, E14, E15, E16, E17, E18). \\
\end{longtable}
}

\begin{center}\rule{0.5\linewidth}{0.5pt}\end{center}

\subsection{Storage / USB}\label{storage-usb}

{\def\LTcaptype{none} % do not increment counter
\begin{longtable}[]{@{}
  >{\raggedright\arraybackslash}p{(\linewidth - 2\tabcolsep) * \real{0.5000}}
  >{\raggedright\arraybackslash}p{(\linewidth - 2\tabcolsep) * \real{0.5000}}@{}}
\toprule\noalign{}
\begin{minipage}[b]{\linewidth}\raggedright
Term
\end{minipage} & \begin{minipage}[b]{\linewidth}\raggedright
Meaning
\end{minipage} \\
\midrule\noalign{}
\endhead
\bottomrule\noalign{}
\endlastfoot
\textbf{LFS / LittleFS} & Flash-friendly file system by ARM; used for
the NAND ``user'' partition. All \texttt{/diags/} data lives here. \\
\textbf{MSC} & USB Mass Storage Class. Mode entered by
\texttt{sentai.usb.drive(1)}; exposes the raw LittleFS block device as
\texttt{/dev/sda} on the host. \\
\textbf{CDC-ACM} & USB Communications Device Class --- Abstract Control
Model. The REPL transport (\texttt{/dev/ttyACM0}). \\
\textbf{CDC-NCM} & USB Communications Device Class --- Network Control
Model. The IP transport (\texttt{enxXXXX} on the host). \\
\textbf{HTTP fast path vs slow path} & See \url{lfs.md} §4.2.
\texttt{/api/raw} reads go through a mutex-protected fast path;
\texttt{/api/ls} always queues to \texttt{lfs\_task} (slow path) to
avoid long \texttt{tcpip\_thread} blocking. \\
\textbf{``Bricked''} & Board visible on USB as Google Coral ID
\texttt{18d1:9307}. Means firmware never started; requires a manual
button press to enter SDP mode. The camera-switch work was careful to
keep USB CDC up before any risky code so the board is never in this
state (see \href{../agent/embeded.md}{embeded.md} §M). \\
\end{longtable}
}

\begin{center}\rule{0.5\linewidth}{0.5pt}\end{center}

\subsection{Measurement methodology}\label{measurement-methodology}

{\def\LTcaptype{none} % do not increment counter
\begin{longtable}[]{@{}
  >{\raggedright\arraybackslash}p{(\linewidth - 2\tabcolsep) * \real{0.5000}}
  >{\raggedright\arraybackslash}p{(\linewidth - 2\tabcolsep) * \real{0.5000}}@{}}
\toprule\noalign{}
\begin{minipage}[b]{\linewidth}\raggedright
Term
\end{minipage} & \begin{minipage}[b]{\linewidth}\raggedright
Meaning
\end{minipage} \\
\midrule\noalign{}
\endhead
\bottomrule\noalign{}
\endlastfoot
\textbf{Warm-up} & Pre-measurement iteration(s) that exercise every
stage on the path but are not recorded. See
\href{../experiments/methodology.md}{experiments/methodology.md}
§3.1. \\
\textbf{Drop-first-sample} & Convention that the first recorded
iteration is discarded from statistics, because it straddles warm-up and
steady state. See \url{statistical_notes.md} §2. \\
\textbf{Sweep} & One block of N iterations with a fixed alternation
pattern. E18 runs three sweeps per session (A fixed cam0, B fixed cam1,
C alternating). \\
\textbf{Per-switch overhead} &
\texttt{C\_total\_mean\ −\ max(A\_total\_mean,\ B\_total\_mean)} in an
E18 head-to-tail session. Table 3 of \url{evaluation.md}. \\
\textbf{Directional asymmetry} &
\texttt{cam1\_total\_mean\ −\ cam0\_total\_mean} within the alternating
sweep of an E16/E18 session. Non-zero before Fix B; within noise after
Fix B. \\
\textbf{Scene snapshot} & JPEG captured from each camera at the start
and end of a session, archived next to the CSVs for offline scene-drift
verification. \\
\textbf{Fault counter} & Persistent-since-boot counter of a
degraded-path event (ISR arm not consumed, drain timeout, grab retry,
grab fatal). Zero-valued post-run counter is a measurement-integrity
check. \\
\end{longtable}
}

\begin{center}\rule{0.5\linewidth}{0.5pt}\end{center}

\subsection{Build / tooling}\label{build-tooling}

{\def\LTcaptype{none} % do not increment counter
\begin{longtable}[]{@{}
  >{\raggedright\arraybackslash}p{(\linewidth - 2\tabcolsep) * \real{0.5000}}
  >{\raggedright\arraybackslash}p{(\linewidth - 2\tabcolsep) * \real{0.5000}}@{}}
\toprule\noalign{}
\begin{minipage}[b]{\linewidth}\raggedright
Term
\end{minipage} & \begin{minipage}[b]{\linewidth}\raggedright
Meaning
\end{minipage} \\
\midrule\noalign{}
\endhead
\bottomrule\noalign{}
\endlastfoot
\textbf{\texttt{flashtool.py}} & \texttt{scripts/flashtool.py} ---
NXP-provided tool; \texttt{-e\ sentai\_runtime} flashes the example in
persistent mode. \\
\textbf{\texttt{\_host\_upload\_repl.py}} &
\texttt{diag/\_host\_upload\_repl.py} --- our REPL-chunked uploader.
Used instead of HTTP \texttt{/api/write/} because the latter hangs on
this firmware. \\
\textbf{\texttt{repl\_run.py}} & Simple REPL driver. Has a stale-prompt
bug on long-running commands; for robust driving we use the inline
\texttt{\_send\_line} pattern from \texttt{\_host\_upload\_repl.py}. \\
\textbf{QSTR} & MicroPython's interned-string pool. Must be regenerated
(with the \texttt{micropython-embed-package} make target) whenever a
\texttt{MP\_QSTR\_*} symbol is added, renamed, or removed. See
\href{../agent/agent.md}{agent/agent.md} §6. \\
\textbf{\texttt{ticks\_ms()}} & \texttt{sentai.rtos.ticks\_ms()} ---
32-bit FreeRTOS tick counter exposed to MicroPython. Wraps every 49.7
days; all deltas are computed with unsigned subtraction so the wrap is
harmless. \\
\textbf{\texttt{\_ticks()}} & Module-local alias for
\texttt{sentai.rtos.ticks\_ms()} used throughout \texttt{diag/}. \\
\end{longtable}
}
